\chapter{On Practice}

Theory is a map. Practice is the road. I have known students who memorised every tag, every elemental resonance, every historical footnote on the Weave's temperament. They could lecture for hours. They could not light a candle without singeing their own sleeve.

This chapter is for those who wish to walk the road, not merely admire the cartography.

\section{The Shape of Learning}

You will fail. Accept this now. The Weave is not impressed by your pedigree, your patron, or your previous successes. It will treat your first spell the same as your thousandth: with polite attention and a receipt.

The difference between a novice and a master is not the absence of failure. It is the speed of recovery.

\subsection{The Three Questions, Restated}

Before every casting, ask yourself three questions. I have written them on the wall of my classroom. I have written them on the inside of my own eyelids.

\begin{enumerate}
    \item \textbf{What do I intend?} Be precise. \emph{``I want to hurt him''} is not precise. \emph{``I want to sear his sword hand so he drops the blade''} is precise. The Weave rewards specificity.
    \item \textbf{Which tags serve this intent?} Name them aloud. \emph{``[BURNING]. [STRIKE]. That is two tags. DV three.''} Speaking the tags is not a ritual. It is a discipline.
    \item \textbf{What is the risk?} Look at your Position. Look at your Fatigue. Look at the witnesses. If the answer is \emph{``I could die,''} consider not casting. Consider running.
\end{enumerate}

A student who answers these questions honestly will still fail. But they will fail \emph{with intention}. And intentional failure is the mother of accidental mastery.

\section{The Practice Regimen}

I recommend the following schedule for novice casters. It is not gentle. It is not accelerated. It is \emph{consistent}.

\subsection{Week One: Cantrips Only}

Cast [LIGHT] until your hand glows in your sleep. Cast [HEAL] on paper cuts. Cast [SILENCE] on your own footsteps. Do not attempt a two- tag spell until you have not singed your sleeve for three consecutive days.

\subsection{Week Two: The Defensive Arts}

Stand in a safe space. Have a friend throw soft objects at you. Cast [DEFEND]. Learn what it feels like to interpose the Weave between yourself and a threat. The friend should throw harder as you improve.

\subsection{Week Three: The Unseen Hand}

Practice [VEIL] in a mirror. Watch your features blur. Say your own name aloud and watch them snap back. Learn the difference between hiding from the world and hiding from yourself.

\subsection{Week Four: The Elemental Chorus}

Find a patch of bare earth. Cast [BURNING] on a dry leaf. Cast [WAVE] on a puddle. Cast [STONE] on a pebble. Cast [WIND] on a feather. Do not move to the next element until the previous one has stopped surprising you.

\subsection{Week Five and Beyond}

You are now ready to fail at more interesting spells. Attempt a [GATE] between two rooms you know well. Attempt a [LEAP] across a gap you could jump without magic. Attempt a [TRANSFORM] on a piece of fruit. The fruit will not thank you.

\section{On Failure, Revisited}

I have failed more spells than most students will ever cast. I keep a list. It is long. It is also instructional.

\subsection{What Failure Teaches}

\begin{itemize}
    \item A miss teaches you that the Weave has opinions. Listen to them.
    \item A partial teaches you that the Weave is willing to negotiate. Learn its terms.
    \item A catastrophic failure teaches you that the Weave has a sense of humour. Do not make it laugh at your expense.
\end{itemize}

The worst failure is not the one that burns your sleeve. It is the one you do not learn from. I have seen students repeat the same mistake for years. They are not practicing. They are rehearsing.

\subsection{The Ritual of the Burnt Sleeve}

When I singe my sleeve—and I still do, occasionally—I cut off the burnt fabric and sew it into a book. That book is now the size of a small pillow. I call it my \emph{Curriculum Vitae}.

You do not need to keep a book of burnt sleeves. But you should keep a record of your failures. Write down what you intended. Write down what you cast. Write down what the Weave said in reply. Read the list before every significant working. It will humble you. Humility is not the enemy of power. It is the container.

\section{The Social Practice}

Magic is not performed in a vacuum. It is performed in the presence of witnesses—allies, enemies, bystanders, and the Weave itself. Your reputation as a caster is not determined by your success rate. It is determined by how you handle your failures.

\subsection{Casting in Company}

When you cast in front of allies, narrate your intention. \emph{``I am veiling our escape. If I fail, we run. Do not wait for me.''} This is not weakness. It is \emph{contract}. The Weave respects contracts.

When you cast in front of enemies, narrate nothing. They do not deserve your grammar.

When you cast in front of bystanders, apologize in advance. \emph{``This may be loud. I am sorry about your windows.''} The Weave appreciates courtesy. So do juries.

\subsection{The Etiquette of Backlash}

When your spell backfires and harms an ally, apologize. Then spend a Boon to help them. The Weave notices reciprocity.

When your spell backfires and harms an enemy, do not apologize. The Weave is not offended by efficiency.

When your spell backfires and harms a bystander, offer compensation. The Weave does not care about silver, but the bystander's lawyer does.

\section{When to Walk Away}

Not every problem requires a magical solution. I have seen students cast [LEAP] to cross a puddle. I have seen others cast [GATE] to avoid a conversation. These are not spells. They are \emph{avoidance}.

\subsection{The Non-Magical Alternative}

Consider, before you cast:

\begin{itemize}
    \item Can you walk around the obstacle?
    \item Can you talk to the guard?
    \item Can you wait for the situation to change?
    \item Can you simply \emph{not}?
\end{itemize}

The Weave is patient. It will wait for you to need it. It does not need you to need it now.

\subsection{The Decision to Stop}

I have stopped casting three times in my life. The first was after a catastrophic failure that killed a horse. The second was after a partial success that blinded me for a day. The third was after a clean success that worked exactly as intended—and I realized I had solved the wrong problem.

Each time, I put down my focus. I walked away. I returned when I remembered why I had started.

You are not a bad caster if you take a break. You are a bad caster if you keep casting when you have forgotten your intention.

\section{The Graduation Test}

At the end of my course, I give students a single spell to cast. They do not know the spell in advance. They do not know the tags. They must construct it in the moment.

The scenario is always the same: a locked door, a burning room, and a single innocent who cannot escape.

The student must choose:

\begin{itemize}
    \item [LEAP] to reach the innocent.
    \item [BURNING] to clear the debris.
    \item [WARD] to protect them from the flames.
    \item [GATE] to open an exit.
    \item Or something else entirely. I have seen [TRANSFORM] turn a door into mist. I have seen [SILENCE] confuse the fire itself. I have seen a student simply open the door—it was not locked—and lead the innocent out.
\end{itemize}

That student passed with highest honors.

\section{A Final Word}

Practice is not repetition. It is \emph{attention}. A student who casts [LIGHT] a thousand times without thinking has learned nothing. A student who casts [LIGHT] ten times, each time noticing the flicker, the color, the Weave's momentary hesitation, has learned more than the thousand- cast student will ever know.

Pay attention. Keep notes. Singe your sleeves. Sew them into a book. One day, that book will be the size of a pillow, and you will look at it and know exactly how much you owe.

The Weave does not forget. Neither should you.

\textit{--- Lysandra}
